\section{Variáveis Aleatórias Discretas e Contínuas}

\begin{def:tipo va}
Uma variável aleatória $X$ é discreta se, e somente se, seu suporte é um conjunto
enumerável. Ela será contínua se, e somente se, exite uma função $f: \mathbb{R} \to \mathbb{R}$
tal que
\[
	F(x) = \int_{-\infty}^x f(u)du
\]
em que $F$ é a função de distribuição acumulada.
\end{def:tipo va}

\paragraph{Exemplos de variáveis aleatórias discretas}
\begin{itemize}
	\item Número de caras obtidas em 10 lançamentos.
	\item Soma dos lados obtidos no lançamento de dados.
	\item Número de clientes em uma loja em um dia.
\end{itemize}
No primeiro exemplo, as realizações possíveis são os inteiros de 0 a 10. Nos dois seguintes,
são todos os naturais. Todos esses são conjunto enumeráveis, cujos pontos isoladamente
concentram massa de probabilidade, isto é, um valor efetivo de probabilidade associado à
ao evento da variável de ter realização naquele ponto.

\paragraph{Exemplos de variáveis aleatórias contínuas}
\begin{itemize}
	\item Tempos de espera em algum serviço.
	\item Peso ou altura de pessoas numa população.
	\item Temperatura ao longo de um ano.
\end{itemize}
Essas variáveis variam de forma contínua num intervalo de valores, de forma que não possível
haver valores que concentrem uma massa de probabilidade, sendo essa totalmente dispersada
pelo intervalo. Em vista disso, a probablidade de que a variável tenha um valor específico
é zero, o que nos leva ao conceito de densidade de probabilidade sobre intervalos de
valores.
